\documentclass{article}
\usepackage[T1]{fontenc}
\usepackage{lmodern}
\usepackage{graphicx}
\usepackage{amsmath,amssymb}
\usepackage[a4paper,margin=2cm]{geometry}
\usepackage[absolute,overlay]{textpos}
\setlength{\TPHorizModule}{1mm}
\setlength{\TPVertModule}{1mm}
\usepackage[indonesian]{babel}
\usepackage{hyperref}
\setlength{\emergencystretch}{2em}
% unit: Halaman 55

\begin{document}

% === Halaman 55 (llm) ===

\section*{2. Acak}

\begin{itemize}
    \item Untuk membuat penyembunyian pesan lebih aman, bit-bit pesan tidak disimpan pada pixel-pixel yang berurutan, namun dipilih secara acak.
    \item Pembangkit bilangan acak-semu (PRNG: \textit{pseudo-random number generator}) digunakan untuk membangkitkan bilangan acak.
    \item Umpan (\textit{seed}) untuk pembangkit bilangan acak berlaku sebagai kunci (\textit{stego-key}).
\end{itemize}

\end{document}
